\paragraph{Data \& models.} This work studies fine-tuning dynamics on a publicly available model (Qwen3-8B under Apache 2.0) and public benchmarks (CoNLL-2003, SQuAD-v1.1, XSum, CodeAlpaca, HumanEval). We do not release new data and we do not train on scraped web content. All model training respects the licence terms of the underlying checkpoint; we do not evaluate on held-out contamination of the benchmarks.

\paragraph{Intended uses.} \tglora{} is intended to reduce the compute and capability cost of adapting LLMs while preserving prior capabilities, including safety-relevant behaviours learned during alignment. The forgetting-topology view gives practitioners a principled tool to \textit{avoid} the unintended erasure of safety properties during downstream fine-tuning, which is an otherwise common failure mode.

\paragraph{Potential misuse.} A methodology that identifies which components carry prior-task behaviour could, in principle, be inverted to identify which components carry \textit{safety} behaviour and to then adapt or overwrite them specifically --- the ``refusal jailbreak via targeted editing'' variant of this concern has been demonstrated in prior ROME-style work. We do not provide recipes or datasets that make this attack easier. All components of our pipeline assume a single known task sequence and measure forgetting against held-out benchmarks, not safety classifiers.

\paragraph{Environmental cost.} Total compute for this paper is estimated in Appendix~\ref{app:repro} using the emissions factor of the grid region of our compute provider. We report this explicitly rather than amortising across groups.

\paragraph{Human annotation.} No new human annotation was collected. All labels come from the original benchmark releases. No study participants were involved.
